\documentclass{beamer}
\usetheme{Madrid}
\usepackage{graphicx}
\usepackage{tikz}
\usepackage{hyperref}
\usepackage{amsmath}
\usepackage{amssymb}

\title{Emotion-Aware Image Caption Generator}
\subtitle{Machine Learning Project Proposal}
% Empty author for footer - this removes names from bottom left
\author{}
% Store the actual author information for the title page only
\newcommand{\authorinfo}{
    \begin{tabular}{rl}
    Team Name: & TeamML\_Newbies \\[0.3em]
    \end{tabular}
    
    \vspace{0.3cm}
    
    \begin{tabular}{rl}
    Anuron Maitro & (2105037) \\
    Md. Asif Ali & (2105038) \\
    Himadri Gobinda Biswas & (2105047)
    \end{tabular}
}

\date{}

\begin{document}

% Slide 1: Title Page
\begin{frame}
    \titlepage
    \vspace{-1cm}
    \begin{center}
        \authorinfo
    \end{center}
\end{frame}

% Slide 2: Project Overview
\begin{frame}{Project Overview}
    \textbf{Project Title:} Emotion-Aware Image Caption Generator \\[0.5em]
    \textbf{Domain:} Computer Vision \& Natural Language Processing \\[1em]
    
    \textbf{Core Idea:}
    \begin{itemize}
        \item Generate descriptive captions for images
        \item Detect emotional expressions in human faces
        \item Integrate emotion information into generated captions
        \item Example: ``A person standing on a road'' $\rightarrow$ ``A person \textit{smiling happily} standing on a road''
    \end{itemize}
    
    \vspace{0.3cm}
    \textbf{Innovation:}
    \begin{itemize}
        \item Goes beyond standard image captioning by adding emotional context
        \item Provides richer, more descriptive image understanding
    \end{itemize}
\end{frame}

% Slide 3: Motivation
% \begin{frame}{Motivation \& Applications}
%     \textbf{Why Emotion-Aware Captions?}
%     \begin{itemize}
%         \item Standard captioning: ``A woman in a park''
%         \item Emotion-aware: ``A woman \textit{looking sad} in a park''
%         \item Emotional context adds significant meaning to descriptions
%     \end{itemize}
    
%     \vspace{0.5cm}
%     \textbf{Real-World Applications:}
%     \begin{itemize}
%         \item \textbf{Accessibility:} Better image descriptions for visually impaired users
%         \item \textbf{Social Media:} Automated caption generation with emotional context
%         \item \textbf{Content Moderation:} Understanding emotional content in images
%         \item \textbf{Healthcare:} Analyzing patient emotional states in medical imaging
%         \item \textbf{Security:} Emotion detection in surveillance systems
%     \end{itemize}
% \end{frame}

% Slide 4: Problem Statement
\begin{frame}{Problem Statement}
    \textbf{Current Limitations of Image Captioning:}
    \begin{itemize}
        \item Existing models describe \textit{what} is in the image
        \item They fail to capture \textit{how} people feel in the image
        \item Emotional state provides critical context for understanding
    \end{itemize}
    
    \vspace{0.5cm}
    \textbf{Our Goal:}
    \begin{itemize}
        \item Build a system that generates captions with emotional awareness
        \item Automatically detect facial expressions and emotions
        \item Intelligently integrate emotion information into captions
        \item Create more human-like, contextually rich descriptions
    \end{itemize}
    
    \vspace{0.5cm}
    \textbf{Challenge:}
    \begin{itemize}
        \item Combine two complex tasks: caption generation + emotion detection
        \item Ensure natural language integration of emotion phrases
    \end{itemize}
\end{frame}

% Slide 5: Proposed Architecture
\begin{frame}{Proposed System Architecture}
    Our system consists of three main components:
    
    \vspace{0.5cm}
    \textbf{Component 1: Base Caption Generation}
    \begin{itemize}
        \item Vision Transformer (ViT) as encoder
        \item GPT-2 as decoder
        \item Fine-tuned on Flickr8k dataset
    \end{itemize}
    
    \vspace{0.3cm}
    \textbf{Component 2: Emotion Detection}
    \begin{itemize}
        \item Pre-trained DeepFace model
        \item Detects 7 emotions: happy, sad, angry, surprise, fear, disgust, neutral
    \end{itemize}
    
    \vspace{0.3cm}
    \textbf{Component 3: Caption Enhancement}
    \begin{itemize}
        \item Custom NLP processing using NLTK
        \item Intelligently inserts emotion phrases into captions
    \end{itemize}
\end{frame}

% % Slide 6: Component 1 - Caption Generation
% \begin{frame}{Component 1: Base Caption Generation}
%     \textbf{Model Architecture:} Vision Transformer + GPT-2
    
%     \vspace{0.5cm}
%     \textbf{Why Vision Transformer (ViT)?}
%     \begin{itemize}
%         \item State-of-the-art performance in image understanding
%         \item Better at capturing global image context than CNNs
%         % \item Pre-trained on ImageNet (14M images)
%         \item Effective encoder for vision-language tasks
%     \end{itemize}
    
%     \vspace{0.5cm}
%     \textbf{Why GPT-2?}
%     \begin{itemize}
%         \item Powerful language generation capabilities
%         % \item Pre-trained on 40GB of text data
%         \item Generates coherent, grammatically correct sentences
%         \item Works well as decoder in vision-to-text tasks
%     \end{itemize}
    
%     \vspace{0.3cm}
%     \textbf{Training Approach:}
%     \begin{itemize}
%         \item Fine-tune the combined ViT-GPT2 model on Flickr8k dataset
%     \end{itemize}
% \end{frame}

% % Slide 7: Component 2 - Emotion Detection
% \begin{frame}{Component 2: Emotion Detection}
%     \textbf{Model:} DeepFace (Pre-trained)
    
%     \vspace{0.5cm}
%     \textbf{Why DeepFace?}
%     \begin{itemize}
%         \item State-of-the-art facial emotion recognition
%         \item Already trained on large facial expression datasets
%         \item No additional training required
%         \item Detects 7 basic emotions with high accuracy
%     \end{itemize}
    
%     \vspace{0.5cm}
%     \textbf{Detected Emotions:}
%     \begin{itemize}
%         \item Happy, Sad, Angry, Surprise, Fear, Disgust, Neutral
%     \end{itemize}
    
%     \vspace{0.5cm}
%     \textbf{Implementation Plan:}
%     \begin{itemize}
%         \item Use pre-trained weights (no training needed)
%         \item Apply to input images to detect facial emotions
%     %     \item Extract dominant emotion with confidence score
%     %     \item Only use emotions with confidence $>$ 30\%
%     \end{itemize}
% \end{frame}

% % Slide 8: Component 3 - Caption Enhancement
% \begin{frame}{Component 3: Caption Enhancement}
%     \textbf{Natural Language Processing with NLTK}
    
%     \vspace{0.5cm}
%     \textbf{Why NLTK?}
%     \begin{itemize}
%         \item Comprehensive toolkit for text processing
%         \item Provides part-of-speech tagging and tokenization
%         \item Allows custom rule-based text manipulation
%         \item Suitable for intelligent phrase insertion
%     \end{itemize}
    
%     \vspace{0.5cm}
%     \textbf{Enhancement Process:}
%     \begin{enumerate}
%         \item Receive base caption from ViT-GPT2 model
%         \item Receive emotion data from DeepFace
%         \item Tokenize caption and perform POS tagging
%         \item Identify person-related nouns (person, man, woman, child)
%         \item Insert emotion phrase after identified person noun
%         \item Return enhanced caption
%     \end{enumerate}
    
%     % \vspace{0.3cm}
%     % \textbf{Example Emotion Phrases:}
%     % \begin{itemize}
%     %     \item Happy $\rightarrow$ ``smiling happily''
%     %     \item Sad $\rightarrow$ ``looking sad''
%     % \end{itemize}
% \end{frame}

% Slide 9: Emotion Phrase Mapping
\begin{frame}{Emotion to Phrase Mapping}
    \textbf{Custom Mapping Strategy:}
    
    \vspace{0.5cm}
    We will create natural descriptive phrases for each emotion:
    
    \vspace{0.5cm}
    \begin{center}
    \begin{tabular}{|l|l|}
    \hline
    \textbf{Detected Emotion} & \textbf{Caption Phrase} \\
    \hline
    Happy & ``smiling happily'' \\
    Sad & ``looking sad'' \\
    Angry & ``looking angry'' \\
    Surprise & ``looking surprised'' \\
    Fear & ``looking worried'' \\
    Disgust & ``looking disgusted'' \\
    Neutral & ``with calm expression'' \\
    \hline
    \end{tabular}
    \end{center}
    
    \vspace{0.5cm}
    These phrases are designed to:
    \begin{itemize}
        \item Sound natural in English sentences
        \item Accurately represent the detected emotion
        % \item Integrate seamlessly into generated captions
    \end{itemize}
\end{frame}

% Slide 10: Dataset
\begin{frame}{Dataset: Flickr8k}
    \textbf{Primary Dataset:} Flickr8k from Kaggle
    
    \vspace{0.5cm}
    \textbf{Dataset Statistics:}
    \begin{itemize}
        \item \textbf{Total Images:} 8,091 images
        \item \textbf{Total Captions:} $\sim$40,000 (5 captions per image)
        \item \textbf{Content:} Diverse real-world scenes with people, objects, activities
    \end{itemize}
    
    \vspace{0.5cm}
    \textbf{Why Flickr8k?}
    \begin{itemize}
        \item Well-established benchmark for image captioning
        \item Contains human-written reference captions
        \item Manageable size for training with limited resources
        \item Diverse image content suitable for testing emotion detection
    \end{itemize}
    
    % \vspace{0.5cm}
    \textbf{Data Split Plan:}
    \begin{itemize}
        \item Training Set: 80\% (6,472 images)
        \item Validation Set: 10\% (809 images)
        \item Test Set: 10\% (810 images)
    \end{itemize}
\end{frame}

% Slide 11: System Pipeline
% \begin{frame}{Complete System Pipeline}
%     \textbf{Step-by-Step Processing:}
    
%     \vspace{0.5cm}
%     \begin{enumerate}
%         \item \textbf{Input:} Receive image from user
%         \vspace{0.2cm}
        
%         \item \textbf{Caption Generation:}
%         \begin{itemize}
%             \item Pass image through ViT encoder
%             \item Generate caption using GPT-2 decoder
%             \item Output: Base caption (e.g., ``A person on a beach'')
%         \end{itemize}
%         \vspace{0.2cm}
        
%         \item \textbf{Emotion Detection:}
%         \begin{itemize}
%             \item Pass image to DeepFace model
%             \item Detect facial emotions if face is present
%             \item Output: Dominant emotion with confidence
%         \end{itemize}
%         \vspace{0.2cm}
        
%         \item \textbf{Caption Enhancement:}
%         \begin{itemize}
%             \item Use NLTK to parse base caption
%             \item Insert emotion phrase at appropriate position
%             \item Output: Enhanced caption (e.g., ``A person \textit{smiling happily} on a beach'')
%         \end{itemize}
%         \vspace{0.2cm}
        
%         \item \textbf{Output:} Return emotion-aware caption to user
%     \end{enumerate}
% \end{frame}

% % Slide 12: Training Strategy
% \begin{frame}{Training Strategy}
%     \textbf{Phase 1: Caption Model Fine-tuning}
%     \begin{itemize}
%         \item Load pre-trained ViT and GPT-2 weights
%         \item Combine into VisionEncoderDecoder architecture
%         \item Fine-tune on Flickr8k training set
%         \item Hyperparameters:
%         \begin{itemize}
%             \item Batch size: 8
%             \item Learning rate: 5e-5
%             \item Epochs: 3-5
%             \item Max caption length: 50 tokens
%         \end{itemize}
%     \end{itemize}
    
%     \vspace{0.5cm}
%     \textbf{Phase 2: Emotion Detection Integration}
%     \begin{itemize}
%         \item Use pre-trained DeepFace (no training needed)
%         \item Apply to test images for emotion extraction
%     \end{itemize}
    
%     \vspace{0.5cm}
%     \textbf{Phase 3: Caption Enhancement}
%     \begin{itemize}
%         \item Implement custom NLTK-based enhancement logic
%         \item Test emotion phrase insertion quality
%     \end{itemize}
% \end{frame}

% % Slide 13: Evaluation Metrics
% \begin{frame}{Evaluation Metrics}
%     We will evaluate our system using both quantitative and qualitative metrics:
    
%     \vspace{0.5cm}
%     \textbf{Caption Quality Metrics:}
%     \begin{itemize}
%         \item \textbf{BLEU Score:} Measures n-gram overlap with reference captions
%         \item \textbf{METEOR:} Evaluates caption quality considering synonyms
%         \item \textbf{ROUGE-L:} Measures longest common subsequence
%         \item \textbf{CIDEr:} Consensus-based image description evaluation
%     \end{itemize}
    
%     \vspace{0.5cm}
%     \textbf{Emotion Detection Metrics:}
%     \begin{itemize}
%         \item Emotion detection success rate
%         \item Confidence score distribution
%     \end{itemize}
    
%     \vspace{0.5cm}
%     \textbf{Human Evaluation:}
%     \begin{itemize}
%         \item Manual assessment of caption naturalness
%         \item Evaluation of emotion relevance and accuracy
%     \end{itemize}
% \end{frame}

% Slide 14: Expected Challenges
% \begin{frame}{Expected Challenges \& Solutions}
%     \textbf{Challenge 1: Emotion Detection Limitations}
%     \begin{itemize}
%         \item Faces may be too small, angled, or occluded
%         \item \textit{Solution:} Set confidence threshold; gracefully handle cases with no face
%     \end{itemize}
    
%     \vspace{0.3cm}
%     \textbf{Challenge 2: Natural Language Integration}
%     \begin{itemize}
%         \item Inserting emotion phrases must sound natural
%         \item \textit{Solution:} Carefully design phrase insertion logic using POS tagging
%     \end{itemize}
    
%     \vspace{0.3cm}
%     \textbf{Challenge 3: Training Resources}
%     \begin{itemize}
%         \item ViT-GPT2 training requires significant GPU resources
%         \item \textit{Solution:} Use Google Colab or Kaggle with GPU acceleration
%     \end{itemize}
    
%     \vspace{0.3cm}
%     \textbf{Challenge 4: Caption Quality}
%     \begin{itemize}
%         \item Generated captions may be generic or inaccurate
%         \item \textit{Solution:} Careful hyperparameter tuning and sufficient training epochs
%     \end{itemize}
% \end{frame}

% Slide 15: Implementation Plan
% \begin{frame}{Implementation Timeline}
%     \textbf{Week 1: Setup \& Data Preparation}
%     \begin{itemize}
%         \item Set up development environment (Kaggle/Colab)
%         \item Download and prepare Flickr8k dataset
%         \item Install required libraries (transformers, deepface, nltk)
%     \end{itemize}
    
%     \vspace{0.3cm}
%     \textbf{Week 2: Base Caption Model}
%     \begin{itemize}
%         \item Implement ViT-GPT2 model architecture
%         \item Fine-tune on Flickr8k training set
%         \item Evaluate caption quality on validation set
%     \end{itemize}
    
%     \vspace{0.3cm}
%     \textbf{Week 3: Emotion Detection \& Integration}
%     \begin{itemize}
%         \item Integrate DeepFace emotion detection
%         \item Implement NLTK-based caption enhancement
%         \item Test complete pipeline on sample images
%     \end{itemize}
    
%     \vspace{0.3cm}
%     \textbf{Week 4: Evaluation \& Refinement}
%     \begin{itemize}
%         \item Comprehensive testing on test set
%         \item Calculate evaluation metrics (BLEU, METEOR, ROUGE)
%         \item Fine-tune and improve results
%     \end{itemize}
% \end{frame}

% % Slide 16: Technology Stack
% \begin{frame}{Technology Stack}
%     \textbf{Programming Language:}
%     \begin{itemize}
%         \item Python 3.8+
%     \end{itemize}
    
%     \vspace{0.3cm}
%     \textbf{Deep Learning Frameworks:}
%     \begin{itemize}
%         \item PyTorch (for model training)
%         \item Hugging Face Transformers (ViT, GPT-2)
%         \item TensorFlow/Keras (for DeepFace)
%     \end{itemize}
    
%     \vspace{0.3cm}
%     \textbf{Key Libraries:}
%     \begin{itemize}
%         \item \textbf{transformers:} Pre-trained ViT and GPT-2 models
%         \item \textbf{deepface:} Facial emotion recognition
%         \item \textbf{nltk:} Natural language processing
%         \item \textbf{opencv:} Image processing
%         \item \textbf{evaluate:} Metric computation (BLEU, METEOR, ROUGE)
%     \end{itemize}
    
%     \vspace{0.3cm}
%     \textbf{Computing Platform:}
%     \begin{itemize}
%         \item Kaggle Notebooks or Google Colab (with GPU)
%     \end{itemize}
% \end{frame}

% % Slide 17: Expected Outcomes
% \begin{frame}{Expected Outcomes}
%     \textbf{Deliverables:}
%     \begin{itemize}
%         \item Fully functional emotion-aware caption generation system
%         \item Fine-tuned ViT-GPT2 model for image captioning
%         \item Integrated emotion detection pipeline
%         \item Well-documented codebase
%         \item Comprehensive evaluation report
%     \end{itemize}
    
%     \vspace{0.5cm}
%     \textbf{Learning Outcomes:}
%     \begin{itemize}
%         \item Understanding of Vision Transformer architecture
%         \item Experience with transfer learning and fine-tuning
%         \item Knowledge of multi-modal AI (vision + language)
%         \item Practical experience with pre-trained models
%         \item Skills in NLP and emotion recognition
%     \end{itemize}
    
%     % \vspace{0.5cm}
%     \textbf{Innovation:}
%     \begin{itemize}
%         \item Novel combination of caption generation and emotion detection
%         \item More descriptive and context-aware image understanding
%     \end{itemize}
% \end{frame}

% % Slide 18: Sample Expected Results
% \begin{frame}{Sample Expected Results}
%     \textbf{Example 1:}
%     \begin{itemize}
%         \item \textbf{Input:} Image of person celebrating
%         \item \textbf{Base Caption:} ``A person jumping in the air''
%         \item \textbf{Detected Emotion:} Happy
%         \item \textbf{Enhanced Caption:} ``A person \textit{smiling happily} jumping in the air''
%     \end{itemize}
    
   
%     \textbf{Example 2:}
%     \begin{itemize}
%         \item \textbf{Input:} Image of person sitting alone
%         \item \textbf{Base Caption:} ``A woman sitting on a bench''
%         \item \textbf{Detected Emotion:} Sad
%         \item \textbf{Enhanced Caption:} ``A woman \textit{looking sad} sitting on a bench''
%     \end{itemize}
    
   
%     \textbf{Example 3:}
%     \begin{itemize}
%         \item \textbf{Input:} Image of landscape (no face)
%         \item \textbf{Base Caption:} ``A mountain with trees''
%         \item \textbf{Detected Emotion:} None
%         \item \textbf{Enhanced Caption:} ``A mountain with trees'' (unchanged)
%     \end{itemize}
% \end{frame}

% New Slide: Project Intent
\begin{frame}{What We Intend to Do}
    \textbf{Core Functional Workflow:}
    \begin{itemize}
        \item \textbf{Universal Base Captioning:} 
        A basic, descriptive caption will be generated for every input image processed by the system.
        
        \item \textbf{Conditional Emotion Analysis:} 
        The system will concurrently analyze the image for emotional expressions using the DeepFace model.
        
        \item \textbf{Selective Enhancement:} 
        An emotion-aware caption will be generated \textit{only} if a specific emotion is successfully detected in the image.
    \end{itemize}

    \vspace{0.5cm}

    \begin{block}{The Processing Logic}
        \centering
        $\text{Image} \rightarrow \text{Base Caption} + (\text{Detected Emotion} \rightarrow \text{Enhanced Caption})$
    \end{block}

    % \vspace{0.3cm}
    % \textbf{Key Objective:}
    % \begin{itemize}
    %     \item To provide contextually rich descriptions when human elements are present, while maintaining standard accuracy for landscape or inanimate scenes[cite: 2, 43].
    % \end{itemize}
\end{frame}

% Slide 18: Sample Expected Results
\begin{frame}{Sample Expected Result}
    \begin{columns}[t]
        % Left Column: Original Image & Base Caption
        \begin{column}{0.45\textwidth}
            \centering
            \includegraphics[width=\textwidth, height=4cm, keepaspectratio]{cycleKid.jpg}
            \vspace{0.2cm}
            \begin{block}{Base Caption}
                \centering \small
                ``A child riding a bicycle on a road'' 
            \end{block}
        \end{column}

        % Right Column: Same Image & Enhanced Caption
        \begin{column}{0.45\textwidth}
            \centering
            \includegraphics[width=\textwidth, height=4cm, keepaspectratio]{cycleKid.jpg}
            \vspace{0.2cm}
            \begin{block}{Enhanced Caption}
                \centering \small
                ``A child \textit{looking worried} riding a bicycle on a road'' 
            \end{block}
        \end{column}
    \end{columns}
    
    % \vspace{0.5cm}
    % \begin{center}
    %     \footnotesize \textit{Note: The system identifies the dominant emotion (e.g., Happy) and integrates it into the base description.} 
    % \end{center}
\end{frame}

% Slide 19: Future Extensions
% \begin{frame}{Potential Future Extensions}
%     \textbf{Short-term Improvements:}
%     \begin{itemize}
%         \item Support for multiple people with different emotions
%         \item More sophisticated emotion phrase generation
%         \item Handling of group emotions (e.g., crowd sentiment)
%     \end{itemize}
    
%     \vspace{0.5cm}
%     \textbf{Long-term Enhancements:}
%     \begin{itemize}
%         \item Real-time video caption generation with emotion tracking
%         \item Multi-lingual emotion-aware captions
%         \item Integration with larger vision-language models (CLIP, BLIP)
%         \item Fine-grained emotion classification (beyond 7 basic emotions)
%         \item Body language and gesture analysis for emotion detection
%     \end{itemize}
    
%     \vspace{0.5cm}
%     \textbf{Applications:}
%     \begin{itemize}
%         \item Mobile app for photo captioning
%         \item Accessibility tool for visual assistance
%         \item Social media content understanding
%     \end{itemize}
% \end{frame}

% Slide 20: References
% \begin{frame}{References}
%     \textbf{Key Papers \& Resources:}
%     \begin{enumerate}
%         \item Dosovitskiy et al., ``An Image is Worth 16x16 Words: Transformers for Image Recognition at Scale'', ICLR 2021
%         \item Radford et al., ``Language Models are Unsupervised Multitask Learners'', OpenAI 2019
%         \item Serengil \& Ozpinar, ``HyperExtended LightFace: A Facial Attribute Analysis Framework'', ICEET 2021 (DeepFace)
%         \item Hodosh et al., ``Framing Image Description as a Ranking Task'', JAIR 2013 (Flickr8k)
%         \item Hugging Face Transformers Documentation
%         \item DeepFace Library Documentation
%     \end{enumerate}
    
%     \vspace{0.5cm}
%     \textbf{Datasets:}
%     \begin{itemize}
%         \item Flickr8k: \url{https://www.kaggle.com/datasets/adityajn105/flickr8k}
%     \end{itemize}
% \end{frame}

% Slide: Existing Works
\begin{frame}{Existing Works}
    \textbf{Related Research:}
    
    \vspace{0.5cm}
    
    \textbf{1. Image Caption Generation using ViT and GPT Architecture}
    \begin{itemize}
        \item IEEE Conference Publication (2024)
        \item \small \url{https://ieeexplore.ieee.org/document/10551257}
    \end{itemize}
    
    \vspace{0.4cm}
    
    \textbf{2. Image Captioning by Incorporating Affective Concepts}
    \begin{itemize}
        \item ScienceDirect (2018)
        \item Uses a very different approach and different dataset from ours
        \item \small \url{https://www.cs.nthu.edu.tw/~lai/pdf/publications/2019/Image_captioning_by_incorporating_affective_concepts_learned_from_both_visual_and_textual_components.pdf}
    \end{itemize}
    
    \vspace{0.4cm}
    
    \textbf{3. Human Emotion Detection Using DeepFace and AI}
    \begin{itemize}
        \item MDPI (2023)
        \item \small \url{https://www.mdpi.com/2673-4591/59/1/37}
    \end{itemize}
\end{frame}

% Slide: Project Progression Plan
\begin{frame}{Project Progression Plan}
    \textbf{25\% Update:}
    \begin{itemize}
        \item Complete environment setup and install all required libraries
        \item Successfully run the project for base image captioning (without emotion detection)
        \item Get a basic caption from image (caption might be gibberish at this point)
    \end{itemize}
    
    \vspace{0.4cm}
    
    \textbf{50\% Update :}
    \begin{itemize}
        \item Implement and integrate emotion detection using DeepFace
        \item Develop NLP-based emotion phrase insertion using NLTK
        \item Define some basic evaluation methodology for emotion-aware captioning(might not comprehensive yet)
        \item Generate emotion-aware captions for sample images(at this point caption with emotion phrase might not be fully syntactically correct)
        \item Try to test  emotion detection accuracy on sample images
    \end{itemize}
    
    \vspace{0.4cm}
\end{frame}
\begin{frame}{Project Progression Plan}    
    \textbf{100\% Update (Final Report):}
    \begin{itemize}
        \item Improving the emotion phrase integration as much as possible
        \item Reduce generation of gibberish captions through extensive testing(produce almost meaningful captions)
        \item Complete comprehensive evaluation on test dataset
        \item Deliver final system with documentation and results analysis
    \end{itemize}
\end{frame}
% Slide: Thank You
\begin{frame}
    \centering
    \Huge Thank You!\\
    % \vspace{1cm}
    % \Large Questions?
\end{frame}

\end{document}